\documentclass[11pt]{article}
\usepackage[margin=0.85in]{geometry}
\usepackage[T1]{fontenc}
\usepackage{amsmath,booktabs,microtype,xurl,hyperref}
\setlength{\emergencystretch}{3em}
\hypersetup{colorlinks=true,urlcolor=blue}
\title{Composition of the Recorded 270.680 ms Controller Call}
\author{Pan Dynamics Collision Avoidance Research}
\date{October 2, 2026}
\begin{document}\maketitle

\section{Finding and measurement boundary}
The requested frame is the first call of the 8 m/s head-on scenario,
at simulation time zero, in the inherited-budget campaign. It uses two
linearization rounds and three numerical conic solves. Numerical optimization
accounts for approximately 26.3 ms, or 9.7\% of the full call. The largest
historical category, 163.545 ms, was not separately instrumented.

New replays locate a large omitted operation: constructing the contracting
terminal core inside the first controller call. Its median measured cost is
106.095 ms when its configuration cache misses, versus 0.107 ms when it hits.
This identifies an important startup bottleneck but does not retroactively
partition every millisecond of the historical residual. The original call's
exact terminal-construction duration remains unknown.

The original 270.680 ms timer surrounds the complete
\path{collisionAvoidanceController} call, including setup and outputs. It
excludes the scenario's later \texttt{ode45} plant replay, rectangle-distance
audit, JSON export and plotting. The campaign had already executed MATLAB
tests, so this was not a fresh MATLAB process; scenario/configuration startup
and interpreter/cache effects can nevertheless remain.

\section{Original recorded decomposition}
These categories are disjoint. CLF construction is subtracted from the
formulation timer before it is reported separately.
\begin{center}
\begin{tabular}{lrr}
\toprule
Recorded operation & Time (ms) & Share (\%) \\
\midrule
Trajectory initialization, predominantly moving-target flow & 15.172 & 5.61 \\
Dynamics/constraints formulation, excluding CLF & 32.560 & 12.03 \\
CLF construction & 28.133 & 10.39 \\
PCBF stages, including the analytic zero-slack case & 5.916 & 2.19 \\
CLF numerical solves & 20.404 & 7.54 \\
Nonlinear prediction-agreement evaluation & 4.950 & 1.83 \\
Unclassified preparation/orchestration/output time & 163.545 & 60.42 \\
\midrule
Full controller call & 270.680 & 100.00 \\
\bottomrule
\end{tabular}
\end{center}
The numerical solver library's internal timers total 24.060 ms; the broader
solve-stage timers include matrix/interface overhead and total 26.306 ms.
The analytic PCBF lower-bound result adds 0.014 ms, bringing the PCBF/CLF
stage total to 26.320 ms. There is no eight-round search in this frame.

\section{Why two rounds and three solves occurred}
The horizon contains eight PCBF prefix holds and sixty completion holds,
at 50 ms per hold: 68 controls and 69 state nodes, spanning 3.4 s.
There is no previous trajectory at time zero, so no inherited safety budget
is available.

\begin{center}
\begin{tabular}{lrr}
\toprule
Operation & Round 1 (ms) & Round 2 (ms) \\
\midrule
Initialization & 15.168 & 0.004 \\
Formulation excluding CLF & 21.554 & 11.006 \\
CLF construction & 15.234 & 12.899 \\
PCBF & 0.014 (analytic) & 5.902 \\
CLF solve & 9.899 & 10.505 \\
Prediction agreement & 2.877 & 2.073 \\
\bottomrule
\end{tabular}
\end{center}
The flow anchor in round 1 satisfies every assembled primary constraint at
zero correction and zero nonnegative slack. Zero is the global lower bound,
so no numerical primary solve is needed. CLF is then optimized numerically
in 14 solver iterations. Its rollout disagrees with the affine prediction:
the body-pose discrepancy measure is 100.463 mm, versus the configured
10 mm tolerance. This measure includes center displacement plus vehicle
reach times heading discrepancy; it is not solely center-position error.
The scaled CLF discrepancy is 0.323042 versus tolerance 0.01, dominating the
maximum normalized agreement ratio of 32.3042.

The controller rebuilds the model along the nonlinear rollout and reduces
the input trust scale from 1 to 0.140754. The new anchor does not satisfy the
complete analytic zero-correction condition, so round 2 solves PCBF
numerically (9 solver iterations), reaching zero optimal safety slack. It
then solves CLF (17 solver iterations). The body-pose discrepancy falls to
2.144 mm, scaled CLF discrepancy to 0.001536, and agreement ratio to
0.214364. The frame terminates immediately. This second round repairs model
accuracy; it is not optional CLF polishing or a failed numerical solve.

\section{What the additional replay establishes}
\path{scripts/runControllerEntryTimingStudy.m} reconstructs the recorded
configuration and initial ego/target state, then injects entry-phase timers
into a temporary renamed copy of the controller. Production code remains
unchanged. Five pairs alternate a terminal cache miss with a hit. Changing
only the constructor's initial radius in a separate, untimed preparation
call evicts its one-entry cache without clearing MATLAB/JIT. Each measured
controller call uses the original configuration and an empty prior state.

\begin{center}
\begin{tabular}{lrr}
\toprule
New measurements (five runs each) & Cache miss & Cache hit \\
\midrule
Median full call (ms) & 184.680 & 79.269 \\
Full-call range (ms) & 183.922--236.499 & 77.320--87.928 \\
Median terminal construction (ms) & 106.095 & 0.107 \\
Median predictive optimization orchestration (ms) & 78.104 & 78.185 \\
\bottomrule
\end{tabular}
\end{center}
Every replay produces exactly the recorded first input to stored numerical
precision, with two rounds and three numerical solves. An unmodified warm
entry call is also compared. Its measured 109.814 ms precedes the paired
series; progressive execution warmup remains visible in the first pair.
The separate instrumented process-warmup call takes 1.462272 s and is not
included in the paired medians. These replays do not replace the historical
270.680 ms observation or establish a hard real-time bound.

A further profiled replay identifies the terminal constructor's work:
one feedback/Lyapunov matrix synthesis, twelve interval RK4 derivative
enclosures (one zero-radius enclosure and eleven radius trials), and eleven
matrix-norm bound calculations. Ten halvings reduce the initial radius
from 0.25 to 0.000244140625, with contraction bound 0.9996028975.
The profile records 48,633 interval-product helper calls and 88,743 outward
rounding helper calls. The matrix interval products use nested scalar loops.
These are terminal-core construction operations, separate from the two
online MPC linearization rounds.

The profiled full call takes 523.727 ms; profiling overhead makes that number
unsuitable for a latency comparison. Inclusive function times overlap and
must not be added together. Call counts and the unprofiled phase timings
support the attribution.

\section{Implication and validation}
The largest identified avoidable startup operation depends on configuration,
not the current encounter: the contracting core can be prepared before
control starts and retained. Existing continuation frames already carry
their terminal object. Reusing or precomputing this mathematical construction
does not require another controller or a changed optimization problem.
No such production change is made in this diagnostic task.

After startup is removed, repeated dynamics/constraint and CLF construction
remain substantial; accelerating the numerical solver alone addresses only
about one tenth of the original call. The full campaign also contains a
262.069 ms noninitial frame, so fixing startup alone would not establish the
requested 100 ms maximum for all frames.

Validation asserts identical configuration, initial state, issued input,
round count and solve count for all measured replays, and unchanged
production entry source. MATLAB R2026a Update 3 runs with
\texttt{-singleCompThread} and the existing native solver. The final Code
Analyzer reports no messages for the new study driver. No controller
algorithm, test suite, scenario or physical limit was modified; no full
campaign rerun is claimed.

The adjacent \path{FRAME_270MS_TIMING_20261002/} directory retains the original
frame, compact summary, replay results, profile function counts/times and
source/artifact hashes. Full profile data and the temporary timed entry are
retained in \path{/home/zai/.cache/collisionAvoidance/frame270-study-20261002/}.
The original controller implementation is commit
\path{9d9f197f6b3dfa38dc1d9fcb666dd51ae971373b}.
\end{document}
